\begin{theorem}[零温四态极限的 Artin--Mazur \texorpdfstring{$\zeta$}{zeta} 函数、迹递推与 primitive 轨道精确主项]\label{thm:xi-zero-temperature-fourstate-artin-mazur-zeta-primitive-asymptotic}
令 \(\Sigma_\infty\) 为由
\[
B_\infty=
\begin{pmatrix}
0&1&1&0\\
0&0&1&0\\
0&0&3&1\\
1&0&0&3
\end{pmatrix}
\]
定义的四态一步 SFT，并记其 Perron 根为 \(\rho:=\rho(B_\infty)\)。再令
\[
a_n:=\#\operatorname{Fix}(\sigma^n|_{\Sigma_\infty})=\Tr(B_\infty^n),
\qquad
\pi_n:=\frac{1}{n}\sum_{d\mid n}\mu(d)\,a_{n/d},
\]
其中 \(\mu\) 为 Möbius 函数。则
\[
\zeta_\infty(z)
:=
\exp\!\left(\sum_{n\ge 1}\frac{a_n}{n}z^n\right)
=
\frac{1}{\det(I-zB_\infty)}
=
\frac{1}{1-6z+9z^2-z^3-z^4}.
\]
此外，序列 \((a_n)_{n\ge 1}\) 满足精确递推
\[
a_{n+4}=6a_{n+3}-9a_{n+2}+a_{n+1}+a_n
\qquad(n\ge 1),
\]
其初值为
\[
a_1=6,\qquad a_2=18,\qquad a_3=57,\qquad a_4=190.
\]
若将 \(B_\infty\) 的特征值按模长排序为
\[
\rho=\lambda_1,
\qquad
\eta:=|\lambda_2|\ge |\lambda_3|\ge |\lambda_4|,
\qquad
\eta<\rho,
\]
则 primitive 轨道数满足
\[
\pi_n=\frac{\rho^n}{n}+O\!\left(\frac{\eta^n}{n}\right).
\]
并且
\[
h_{\mathrm{top}}(\Sigma_\infty)=\log\rho.
\]
等价地，在推论 \ref{cor:real-input-40-ground-entropy} 所采用的原始系统步长上，零温地基每步熵密度为 \(\tfrac12\log\rho\)。
\end{theorem}
\begin{proof}
零温四态极限段已给出
\[
\det(I-zB_\infty)=1-6z+9z^2-z^3-z^4.
\]
对任意一步 SFT，Artin--Mazur \(\zeta\)-函数满足
\[
\zeta(z)=\det(I-zB)^{-1},
\qquad
\#\operatorname{Fix}(\sigma^n)=\Tr(B^n),
\]
故第一式立即成立。

\(B_\infty\) 的特征多项式为
\[
\chi(\lambda)=\lambda^4-6\lambda^3+9\lambda^2-\lambda-1.
\]
由 Cayley--Hamilton 定理，
\[
B_\infty^{n+4}=6B_\infty^{n+3}-9B_\infty^{n+2}+B_\infty^{n+1}+B_\infty^{n}
\qquad(n\ge 1),
\]
两侧取迹即得递推。初值由直接矩阵乘法给出
\[
\Tr(B_\infty)=6,\qquad
\Tr(B_\infty^2)=18,\qquad
\Tr(B_\infty^3)=57,\qquad
\Tr(B_\infty^4)=190.
\]

再看 primitive 轨道主项。对 \(\chi\) 直接求根可知其四个根两两不同，且按模长依次为
\[
\rho\approx 3.596154676009,\qquad
\eta\approx 2.182644323596,\qquad
|\lambda_3|\approx 0.484278410709,\qquad
|\lambda_4|\approx 0.263077410313.
\]
特别地，
\[
\eta>\sqrt{\rho}.
\]
因此由谱分解，
\[
a_n=\lambda_1^n+\lambda_2^n+\lambda_3^n+\lambda_4^n
=
\rho^n+O(\eta^n).
\]
再由 Möbius 反演公式
\[
n\pi_n=\sum_{d\mid n}\mu(d)\,a_{n/d}
=
a_n+\sum_{\substack{d\mid n\\ d\ge 2}}\mu(d)\,a_{n/d},
\]
其中主项 \(a_n=\rho^n+O(\eta^n)\)。对任意真约数 \(d\ge 2\)，有 \(n/d\le n/2\)，故
\[
a_{n/d}=O(\rho^{n/d})=O(\rho^{n/2}).
\]
又因 \(\eta>\sqrt{\rho}\)，故 \(\rho^{n/2}=O(\eta^n)\)；再结合约数函数 \(\tau(n)=n^{o(1)}\) 的次指数增长，可得
\[
\sum_{\substack{d\mid n\\ d\ge 2}}\mu(d)\,a_{n/d}=O(\eta^n).
\]
于是
\[
n\pi_n=\rho^n+O(\eta^n),
\]
即
\[
\pi_n=\frac{\rho^n}{n}+O\!\left(\frac{\eta^n}{n}\right).
\]

最后，边移位 \(\Sigma_\infty\) 的拓扑熵等于 \(\log\rho(B_\infty)\)；与推论 \ref{cor:real-input-40-ground-entropy} 比较，即得原始系统步长上的熵密度为 \(\tfrac12\log\rho\)。证毕。
\end{proof}
